\documentclass[12pt,a4paper]{article}

% --- Scientific Reports (Springer Nature) packages ---
\usepackage[utf8]{inputenc}
\usepackage[T1]{fontenc}
\usepackage{times}
\usepackage[margin=2.5cm]{geometry}
\usepackage{amsmath,amssymb,amsfonts}
\usepackage{graphicx}
\usepackage{booktabs}
\usepackage{multirow}
\usepackage{hyperref}
\usepackage{float}
\usepackage{algorithm}
\usepackage{algorithmic}
\usepackage{xcolor}
\usepackage{enumitem}
\usepackage{array}
\usepackage{caption}
\usepackage{setspace}
\usepackage{natbib}
\usepackage{lineno}

% Scientific Reports uses line numbers
\linenumbers
\onehalfspacing

\title{\textbf{TH-DAT: Trimester-Aware Hierarchical Domain-Attention Transformer for Early Prediction of Antenatal Depression Using Multi-Cohort Clinical Data}}

\author{
Saivenkat~H\textsuperscript{1}, Lavanya\textsuperscript{1}, Adithya\textsuperscript{1}, Shapper\textsuperscript{1}, Beebi~Naseeba\textsuperscript{2,*}, Rishitha~Muddana\textsuperscript{2,*} \\[6pt]
\textsuperscript{1}School of Computer Science and Engineering, VIT-AP University, Amaravati, Andhra Pradesh, India \\
\textsuperscript{2}Assistant Professor, School of Computer Science and Engineering, VIT-AP University, Amaravati, Andhra Pradesh, India \\[6pt]
\textsuperscript{*}Corresponding guides and supervisors. \\
Email: hsaivenkat1@gmail.com
}

\date{}

\begin{document}

\maketitle

% =====================================================================
% ABSTRACT
% =====================================================================
\begin{abstract}
\noindent
Antenatal depression remains a critically underdiagnosed condition affecting 10--20\% of pregnant women worldwide, contributing to adverse maternal and neonatal outcomes including preterm birth, low birth weight, impaired maternal-infant bonding, and elevated postnatal depression risk. Current machine learning approaches for depression risk prediction treat clinical features as flat, unstructured inputs, ignoring both the clinically meaningful domain groupings (demographic, obstetric, psychosocial) and the temporal progression of risk factors across pregnancy trimesters. We propose TH-DAT (Trimester-Aware Hierarchical Domain-Attention Transformer), a novel deep learning architecture that introduces four synergistic innovations: (1)~domain-grouped feature tokenization with domain-aware positional encoding aligned with clinical knowledge, (2)~a global transformer encoder enabling both intra-domain and cross-domain feature interactions, (3)~trimester-sequential cross-attention conditioning predictions on pregnancy stage, and (4)~gated domain fusion with entropy regularization providing patient-level interpretability. TH-DAT employs a two-phase training strategy combining masked-feature self-supervised pretraining with focal-contrastive supervised fine-tuning, and applies SMOTE for addressing class imbalance. Evaluated on two geographically independent antenatal datasets---PERI\_DEP (N=14,008, Pakistan) and Uganda EPDS (N=14,325)---TH-DAT achieves state-of-the-art AUC-ROC scores of \textbf{0.9769} and \textbf{0.9610} respectively, outperforming all 11 baseline models including Random Forest (0.9670), XGBoost (0.9072), TabTransformer (0.9527), SAINT (0.8543), and text transformers (BERT, GPT-2, RoBERTa, T5). A systematic ablation study across six configurations confirms the contribution of each architectural component, with masked pretraining ($-$0.98\% AUC) and domain grouping ($-$0.95\% AUC) providing the largest performance gains. SHAP analysis independently validates the balanced domain gate weights (Demographic: 34.3\%, Obstetric: 32.8\%, Psychosocial: 32.9\%), with less than 1\% discrepancy between model-agnostic SHAP attributions and the model's built-in gate weights. These gate weights offer clinically actionable interpretability absent from all competing approaches, enabling personalized intervention recommendations based on each patient's dominant risk domain.
\end{abstract}

\noindent\textbf{Keywords:} Antenatal depression, transformer architecture, tabular data classification, attention mechanism, clinical decision support, interpretable artificial intelligence, maternal mental health, self-supervised pretraining, domain-aware learning, SMOTE, SHAP, multi-cohort validation

\newpage

% =====================================================================
% INTRODUCTION
% =====================================================================
\section{Introduction}

Depression during pregnancy---termed antenatal depression (AD)---represents the most prevalent psychiatric complication of the perinatal period, affecting an estimated 10 to 20 percent of women globally\cite{woody2017}. The World Health Organization recognizes maternal mental health disorders as a leading cause of disease burden among women of reproductive age, with disproportionate impact in low- and middle-income countries (LMICs) where up to 25\% of pregnant women experience clinically significant depressive symptoms\cite{fisher2012}. Untreated antenatal depression carries significant consequences for both mother and offspring, including elevated risk of preterm birth (OR~1.37), intrauterine growth restriction, low birth weight (OR~1.28), compromised maternal-infant bonding, inadequate prenatal care utilization, substance use during pregnancy, and a 3--5 fold increased probability of postnatal depression\cite{grote2010, howard2014}.

Despite its clinical significance and well-documented adverse outcomes, systematic screening for antenatal depression remains inconsistent across healthcare systems globally. Standard screening protocols rely on self-report instruments such as the Edinburgh Postnatal Depression Scale (EPDS)\cite{cox1987} and the Patient Health Questionnaire (PHQ-9)\cite{kroenke2001}, which require trained personnel for administration and interpretation. In resource-constrained settings---where the burden of antenatal depression is paradoxically highest---routine screening is frequently omitted due to workforce limitations, competing clinical priorities during prenatal visits, stigma surrounding mental health, and lack of mental health infrastructure and referral pathways\cite{rahman2008}.

An automated predictive system that leverages routinely collected clinical data---demographic background, obstetric history, and psychosocial indicators---could enable proactive identification of at-risk women during standard antenatal care visits, without requiring dedicated screening instruments or specialized mental health training. Such a system must satisfy three critical requirements: (1)~high predictive accuracy to minimize both missed cases (false negatives) and unnecessary referrals (false positives), (2)~clinical interpretability enabling practitioners to understand and act upon predictions, and (3)~cross-population generalizability ensuring reliable performance across diverse healthcare contexts.

\subsection{Research Challenges}

The development of accurate, generalizable, and interpretable predictive models for antenatal depression faces several interconnected technical and clinical challenges:

\begin{enumerate}[label=\textbf{C\arabic*:}]
    \item \textbf{Domain Structure Ignorance:} Clinical features naturally cluster into demographic (age, education, employment), obstetric (gestational age, parity, complications), and psychosocial (physical health perception, body image, family relationships) domains. Existing ML models treat these as flat, exchangeable vectors, discarding the intra-domain correlations (e.g., education correlates with employment within demographics) and inter-domain interactions (e.g., obstetric complications exacerbate psychosocial stress) that are critical for accurate risk assessment.
    
    \item \textbf{Temporal Insensitivity:} Depression risk profiles evolve across pregnancy trimesters. First-trimester risk factors (unplanned pregnancy, morning sickness, adjustment difficulties) differ substantially from third-trimester factors (delivery anxiety, financial preparation stress, physical discomfort). No existing computational model captures this temporal dimension of pregnancy-related depression risk.
    
    \item \textbf{Class Imbalance:} Depression prevalence in antenatal populations typically ranges from 15--35\%, creating significant class imbalance. Models trained with standard cross-entropy loss achieve deceptively high accuracy by preferentially predicting the majority class (non-depressed), resulting in poor sensitivity---precisely the patients who need clinical attention are systematically missed.
    
    \item \textbf{Interpretability Deficit:} Clinical adoption of AI-based screening tools requires transparent, actionable predictions\cite{amann2020}. Clinicians need to understand \textit{why} a patient is flagged for depression risk in terms of clinically meaningful factors, not merely receive an opaque probability score. Post-hoc interpretability methods such as SHAP\cite{lundberg2017} and LIME\cite{ribeiro2016} provide feature-level explanations but are computationally expensive and may produce inconsistent explanations.
    
    \item \textbf{Generalizability Concerns:} Models developed and validated on single datasets from specific populations may not generalize across cohorts with different cultural, socioeconomic, and healthcare contexts. Cross-cohort validation using geographically and culturally diverse datasets is essential yet rarely performed in perinatal depression prediction research.
\end{enumerate}

\subsection{Research Objectives}

This work addresses the five challenges identified above through the following objectives:

\begin{enumerate}
    \item Design a novel transformer architecture that embeds clinical domain knowledge and trimester awareness directly into the attention mechanism rather than treating features as an unstructured flat vector.
    \item Develop an interpretable gated fusion mechanism that produces per-patient domain contribution weights without relying on post-hoc explanation methods.
    \item Implement a two-phase training strategy that synergistically addresses both representation learning (via masked pretraining) and class imbalance (via focal loss and SMOTE).
    \item Establish the most comprehensive benchmark in antenatal depression prediction, comparing 12 models spanning four categories (classical ML, deep learning, tabular transformers, text transformers) on two international datasets.
    \item Validate model generalizability through multi-cohort evaluation on geographically independent datasets from Pakistan and Uganda.
\end{enumerate}

\subsection{Contributions}

The principal contributions of this work are:

\begin{itemize}
    \item \textbf{Novel Architecture:} TH-DAT is the first transformer architecture specifically designed for antenatal clinical tabular data, incorporating domain-grouped tokenization, trimester cross-attention, and gated interpretable fusion---capabilities absent from all prior approaches.
    \item \textbf{State-of-the-Art Performance:} TH-DAT achieves the highest AUC-ROC (0.9769) among all 12 models evaluated, outperforming Random Forest (0.9670, +0.99\%), TabTransformer (0.9527, +2.42\%), and BERT (0.6760, +30.09\%).
    \item \textbf{Two-Phase Training with SMOTE:} Combined masked-feature pretraining, focal-contrastive fine-tuning, and SMOTE oversampling addressing representation learning, loss sensitivity, and data-level class imbalance simultaneously.
    \item \textbf{Comprehensive Benchmark:} Comparison of 12 models across four categories on two international datasets---the most extensive evaluation in perinatal depression prediction.
    \item \textbf{Systematic Ablation:} Six-configuration ablation study quantifying each architectural component's contribution to overall performance.
    \item \textbf{Clinical Interpretability:} Per-patient domain gate weights providing actionable clinical insights, independently validated by SHAP analysis with $<$1\% domain-level discrepancy.
    \item \textbf{Multi-Cohort Validation:} Consistent performance across Pakistani (PERI\_DEP) and Ugandan (EPDS) populations with different cultural contexts and depression prevalence rates.
\end{itemize}

The remainder of this paper is organized as follows. Section~2 provides a comprehensive literature survey. Section~3 formulates the problem mathematically. Section~4 presents the proposed TH-DAT architecture in detail. Section~5 describes datasets and preprocessing including SMOTE. Section~6 presents the experimental setup. Sections~7--8 report primary results and extended metrics. Section~9 analyzes ROC and PR curves. Section~10 presents confusion matrix analysis. Section~11 provides comparative analysis. Section~12 reports ablation study results. Section~13 analyzes robustness, complexity, and SHAP interpretability. Section~14 discusses findings and clinical implications. Section~15 concludes the paper.

% =====================================================================
% LITERATURE SURVEY
% =====================================================================
\section{Literature Survey}

The application of machine learning (ML) and deep learning (DL) to perinatal depression prediction has grown substantially in recent years. This section reviews 18 representative studies organized by methodology, followed by analysis of tabular transformer architectures, text-based approaches, and explainable AI methods. Table~\ref{tab:litsurvey} provides a comprehensive comparison of all existing approaches.

\subsection{Machine Learning Classifiers for Perinatal Depression}

Krishnamurti et al.\cite{krishnamurti2024} applied machine learning classifiers to patient-reported pregnancy data for predicting first-time depression onset during pregnancy, achieving an AUC of 0.89. Their work demonstrated the predictive value of self-reported clinical information but relied on traditional flat-feature approaches without domain-specific processing. Huang et al.\cite{huang2024bias} specifically addressed model bias in prenatal depression prediction using ML with bias assessment techniques, though their best model achieved only 0.61 AUC, highlighting the fundamental challenge of simultaneously ensuring fairness and predictive accuracy in clinical ML systems.

Hu et al.\cite{hu2025} developed and applied an ML-based antenatal depression prediction model using clinical maternal records, achieving 0.710 AUC. Their relatively lower performance underscores the difficulty of prediction from routine clinical data alone. Chen et al.\cite{chen2025} took a novel approach by combining elastic net regression with ML, leveraging biochemical markers including remnant cholesterol for peripartum depression prediction (0.875 AUC), demonstrating the potential of multimodal feature integration.

\subsection{Multiple Model Development and Validation}

Qi et al.\cite{qi2025} systematically developed and validated multiple ML models for postpartum depression prediction in women using maternal clinical data, reaching 0.909 AUC with their best-performing model. Wang et al.\cite{wang2025} pursued interpretable ML for predicting perinatal depression in Chinese women during mid- and late pregnancy, achieving 0.917 AUC with both development and external validation cohorts---one of the higher-performing prior studies. Han et al.\cite{han2023} developed a comprehensive psychological tendency prediction model specifically for pregnant women based on standardized questionnaires, reporting 92.3\% accuracy through feature engineering of questionnaire responses.

\subsection{Ensemble and High-Risk Prediction}

Khadidos et al.\cite{khadidos2024} proposed an ensemble ML framework for predicting maternal health risk during pregnancy, achieving notably high accuracy of 98.4\%, though this addressed general maternal health risk rather than depression specifically. Pi et al.\cite{pi2025} applied various ML algorithms to high-risk pregnancy prediction based on clinical features, achieving 94.6\% accuracy.

\subsection{Novel Data Modalities}

Wiley et al.\cite{wiley2025} explored a fundamentally different approach through voice biomarkers for perinatal depression detection in a cross-sectional nationwide pilot study, achieving 0.78 AUC using speech recordings---demonstrating the potential of multimodal data sources beyond structured clinical records. Zhang et al.\cite{zhang2024speech} similarly combined LightGBM with speech features for depression and anxiety screening in pregnant women using free conversational speech in naturalistic conditions, achieving 0.79 AUC.

\subsection{Risk Factor Identification and Interpretability}

Sibbald et al.\cite{sibbald2025} focused on identifying prenatal risk factors of postpartum depression using LASSO-regularized ML with prenatal psychological, biological, and social data (0.72 AUC). Lin and Zhou\cite{lin2025} proposed an ensemble neural network method for predicting postpartum depression, achieving 90.0\% accuracy. Huang et al.\cite{huang2025explain} applied explainable ML algorithms to postpartum depression risk prediction using clinical and demographic maternal data (0.91 AUC). Zhang et al.\cite{zhang2025interp} developed an interpretable ML model for predicting postpartum depression through a retrospective study design (0.89 AUC). Kim\cite{kim2025} conducted predictive analysis of postpartum depression using ML on postpartum clinical datasets (0.88 AUC).

\subsection{Longitudinal and Large-Scale Studies}

Wang et al.\cite{wang2024traj} studied the trajectory of postpartum depression in Chinese women using a three-wave longitudinal follow-up research design with ML (0.84 AUC), providing insights into temporal evolution of depression symptoms. Malani et al.\cite{malani2026} developed an ML prediction model for postpartum psychiatric admission using gradient-boosted trees and elastic net on the large-scale Born in Queensland cohort comprising 205,495 pregnancies, achieving 0.79 AUC.

\begin{table}[!htbp]
\caption{Comparison of Existing Machine Learning-Based Studies for Antenatal and Perinatal Depression Prediction}
\label{tab:litsurvey}
\centering
\small
\renewcommand{\arraystretch}{1.25}
\begin{tabular}{|c|p{6.5cm}|p{2.2cm}|p{2.3cm}|p{2.5cm}|c|}
\hline
\textbf{S.No.} & \textbf{Paper Title} & \textbf{Authors} & \textbf{Method} & \textbf{Dataset} & \textbf{Accuracy} \\
\hline
1 & Predicting First Time Depression Onset in Pregnancy: Applying Machine Learning Methods to Patient-Reported Data & Krishnamurti et al. (2024) & Machine Learning Classifiers & Patient-reported pregnancy data & 0.89 AUC \\
\hline
2 & Predicting Prenatal Depression and Assessing Model Bias Using Machine Learning Models & Huang et al. (2024) & ML + Bias Assessment & Prenatal clinical records & 0.61 AUC \\
\hline
3 & Development and Application of a Machine Learning-Based Antenatal Depression Prediction Model & Hu et al. (2025) & Machine Learning Prediction Model & Clinical maternal records & 0.710 AUC \\
\hline
4 & Predicting Peripartum Depression Using Elastic Net Regression and Machine Learning: The Role of Remnant Cholesterol & Chen et al. (2025) & Elastic Net + ML & Clinical and biochemical data & 0.875 AUC \\
\hline
5 & Prediction of Postpartum Depression in Women: Development and Validation of Multiple Machine Learning Models & Qi et al. (2025) & Multiple ML Models & Maternal clinical data & 0.909 AUC \\
\hline
6 & Development and External Validation of an Interpretable Machine Learning Model for Predicting Perinatal Depression in Chinese Women During Mid- and Late Pregnancy & Wang et al. (2025) & Explainable Machine Learning & Chinese maternal cohort & 0.917 AUC \\
\hline
7 & A Comprehensive Psychological Tendency Prediction Model for Pregnant Women Based on Questionnaires & Han et al. (2023) & Psychological Prediction Model & Questionnaire responses & 92.3\% \\
\hline
8 & Ensemble Machine Learning Framework for Predicting Maternal Health Risk During Pregnancy & Khadidos et al. (2024) & Ensemble Machine Learning & Maternal health risk dataset & 98.4\% \\
\hline
9 & Prediction of High-Risk Pregnancy Based on Machine Learning Algorithms & Pi et al. (2025) & Machine Learning Algorithms & High-risk pregnancy dataset & 94.6\% \\
\hline
10 & Voice Biomarkers of Perinatal Depression: Cross-Sectional Nationwide Pilot Study Report & Wiley et al. (2025) & Voice Biomarker Analysis & Speech recordings & 0.78 AUC \\
\hline
11 & Identifying Prenatal Risk Factors of Postpartum Depression with Machine Learning & Sibbald et al. (2025) & LASSO-Regularized ML & Prenatal psychological, biological and social data & 0.72 AUC \\
\hline
12 & A Method for Predicting Postpartum Depression via an Ensemble Neural Network Model & Lin and Zhou (2025) & Ensemble Neural Network & Postpartum depression clinical dataset & 90.0\% \\
\hline
13 & Postpartum Depression Risk Prediction Using Explainable Machine Learning Algorithms & Huang et al. (2025) & Explainable Machine Learning & Clinical and demographic maternal data & 0.91 AUC \\
\hline
14 & Interpretable Machine Learning Model for Predicting Postpartum Depression: Retrospective Study & Zhang et al. (2025) & Interpretable Machine Learning & Retrospective clinical dataset & 0.89 AUC \\
\hline
15 & Predictive Analysis of Postpartum Depression Using Machine Learning & Kim (2025) & Machine Learning & Postpartum clinical dataset & 0.88 AUC \\
\hline
16 & Depression and Anxiety Screening for Pregnant Women via Free Conversational Speech in Naturalistic Condition & Zhang et al. (2024) & LightGBM + Speech Features & Conversational speech recordings & 0.79 AUC \\
\hline
17 & Trajectory on Postpartum Depression of Chinese Women and the Risk Prediction Models: A Machine-Learning Based Three-Wave Follow-Up Research & Wang et al. (2024) & Machine Learning & Longitudinal postpartum cohort & 0.84 AUC \\
\hline
18 & Developing a Machine Learning Prediction Model for Postpartum Psychiatric Admission: Findings from the Born in Queensland Study & Malani et al. (2026) & Gradient-Boosted Trees, Elastic Net & Born in Queensland cohort (205,495 pregnancies) & 0.79 (AUC) \\
\hline
\textbf{19} & \textbf{TH-DAT: Trimester-Aware Hierarchical Domain-Attention Transformer for Early Prediction of Antenatal Depression (Ours)} & \textbf{Proposed (2026)} & \textbf{Domain-Aware Transformer + SMOTE + SHAP} & \textbf{PERI\_DEP + Uganda EPDS (28,333 records)} & \textbf{0.977 AUC} \\
\hline
\end{tabular}
\end{table}

\subsection{Transformer Architectures for Tabular Data}

Huang et al.\cite{huang2020tabtf} introduced TabTransformer, the first application of multi-head self-attention to tabular data, embedding categorical features through learned lookup tables while concatenating numerical features directly. Gorishniy et al.\cite{gorishniy2021} proposed FT-Transformer, tokenizing all features into a shared embedding space with a [CLS] classification token. Somepalli et al.\cite{somepalli2021} introduced SAINT with intersample attention combined with contrastive pretraining. Despite these advances, all existing tabular transformers lack domain-specific feature processing, temporal conditioning, and built-in interpretability.

\subsection{Language Models for Structured Data}

Row serialization approaches convert tabular records to natural language text strings and process them with pretrained language models such as BERT\cite{devlin2019}, GPT-2\cite{radford2019}, RoBERTa\cite{liu2019}, and T5\cite{raffel2020}. Dinh et al.\cite{dinh2022} and Hegselmann et al.\cite{hegselmann2023} demonstrated that while large language models can process serialized tabular data, they fundamentally struggle with numerical precision and cross-feature interactions.

\subsection{Explainable AI in Clinical Settings}

SHAP\cite{lundberg2017} provides theoretically grounded feature attributions based on cooperative game theory. LIME\cite{ribeiro2016} approximates model behavior locally through interpretable surrogate models. Jain and Wallace\cite{jain2019} demonstrated that standard attention weights do not reliably indicate feature importance. Our gated domain fusion addresses this by producing explicit, constrained domain-level weights.

\subsection{Research Gap Analysis}

Across all 18 studies in Table~\ref{tab:litsurvey}, no prior work simultaneously addresses: (1)~domain-aware feature processing, (2)~trimester-based temporal conditioning, (3)~gated interpretable fusion, (4)~self-supervised masked pretraining for tabular clinical data, and (5)~multi-cohort validation. TH-DAT is the first architecture to address all five desiderata simultaneously.

% =====================================================================
% PROBLEM FORMULATION
% =====================================================================
\section{Problem Formulation}

\subsection{System Model and Notation}

Let $\mathcal{D} = \{(\mathbf{x}_i, t_i, y_i)\}_{i=1}^{M}$ denote a clinical dataset of $M$ antenatal patient records, where:
\begin{itemize}
    \item $\mathbf{x}_i = [x_{i,1}, x_{i,2}, \ldots, x_{i,N}] \in \mathbb{R}^N$ is the feature vector containing $N$ clinical measurements
    \item $t_i \in \{0, 1, 2\}$ is the trimester indicator derived from gestational age
    \item $y_i \in \{0, 1\}$ is the binary depression label
\end{itemize}

The goal is to learn a mapping $f_\theta: \mathbb{R}^N \times \{0,1,2\} \rightarrow [0,1]$ predicting depression probability.

\subsection{Domain Partitioning}
\begin{equation}
\mathcal{I}_{\text{demo}} \cup \mathcal{I}_{\text{obst}} \cup \mathcal{I}_{\text{psyc}} = \{1, \ldots, N\}, \quad \mathcal{I}_k \cap \mathcal{I}_j = \emptyset \ \forall k \neq j
\end{equation}

For PERI\_DEP ($N=19$): $|\mathcal{I}_{\text{demo}}| = 6$, $|\mathcal{I}_{\text{obst}}| = 10$, $|\mathcal{I}_{\text{psyc}}| = 3$. For Uganda EPDS ($N=13$): $|\mathcal{I}_{\text{demo}}| = 4$, $|\mathcal{I}_{\text{obst}}| = 5$, $|\mathcal{I}_{\text{psyc}}| = 4$.

\subsection{Trimester Derivation}
\begin{equation}
t = \begin{cases} 0 & \text{if gestational age } g \leq 13 \text{ weeks (T1)} \\ 1 & \text{if } 14 \leq g \leq 26 \text{ weeks (T2)} \\ 2 & \text{if } g \geq 27 \text{ weeks (T3)} \end{cases}
\end{equation}

\subsection{Optimization Objective}
\begin{equation}
\theta^* = \arg\min_\theta \frac{1}{M} \sum_{i=1}^{M} \left[ \mathcal{L}_{\text{focal}}(\hat{y}_i, y_i) + \lambda_1 \mathcal{L}_{\text{con}}(\mathbf{z}_i, y_i) + \lambda_2 \mathcal{L}_{\text{ent}}(\mathbf{g}_i) \right]
\end{equation}
where $\lambda_1 = 0.03$, $\lambda_2 = 0.01$.

% =====================================================================
% PROPOSED METHOD
% =====================================================================
\section{Proposed TH-DAT Architecture}

\subsection{Architectural Overview}

TH-DAT processes clinical tabular data through five hierarchical components (Fig.~\ref{fig:arch}), parameterized by approximately 330,916 trainable parameters.

\begin{figure}[htbp]
\centering
\includegraphics[width=\textwidth]{thdat_architecture.jpg}
\caption{Complete TH-DAT architecture. Patient records are partitioned into three clinical domains---Demographic (blue), Obstetric (orange), Psychosocial (green)---and tokenized with domain-aware positional embeddings (C1). A 4-layer global transformer encoder with trimester conditioning processes all tokens (C2). Domain-specific attention pooling compresses tokens into domain summaries (C3). Trimester cross-attention produces a pregnancy-stage-conditioned representation (C4). Gated domain fusion generates interpretable per-patient weights $g_1, g_2, g_3$ (C5). A skip connection preserves raw feature information.}
\label{fig:arch}
\end{figure}

\subsection{C1: Domain-Grouped Feature Tokenizer}

Each scalar feature $x_i \in \mathbb{R}$ is projected into $d$-dimensional embedding space:
\begin{equation}
\mathbf{e}_i = \text{LayerNorm}(\text{GELU}(W_i x_i + b_i)), \quad W_i \in \mathbb{R}^{d \times 1}
\end{equation}

TH-DAT uses per-feature projections ($N$ separate linear layers). Tokens are augmented with dual positional information:
\begin{equation}
\mathbf{h}_i = \mathbf{e}_i + \mathbf{p}_i + \mathbf{d}_{D(i)}
\end{equation}
where $\mathbf{p}_i \in \mathbb{R}^d$ is positional embedding and $\mathbf{d}_{D(i)} \in \mathbb{R}^d$ is domain embedding.

\subsection{C2: Global Transformer Encoder with Trimester Conditioning}

Trimester conditioning is injected additively:
\begin{equation}
\mathbf{h}_i' = \mathbf{h}_i + \alpha \cdot \text{Embedding}_{\text{tri}}(t), \quad \alpha = 0.1
\end{equation}

The transformer encoder processes:
\begin{equation}
\mathbf{H}'' = \text{TransformerEncoder}(\mathbf{H}'; L=4, H=8, d_{ff}=256)
\end{equation}

Multi-head self-attention:
\begin{equation}
\text{MHA}(Q, K, V) = \text{Concat}(\text{head}_1, \ldots, \text{head}_H) W^O
\end{equation}
\begin{equation}
\text{head}_j = \text{Softmax}\left(\frac{(QW_j^Q)(KW_j^K)^\top}{\sqrt{d_k}}\right)(VW_j^V)
\end{equation}

Position-wise FFN: $\text{FFN}(\mathbf{x}) = W_2 \cdot \text{GELU}(W_1 \mathbf{x} + b_1) + b_2$.

\subsection{C3: Domain-Aware Attention Pooling}

For each domain $k$, a learnable query $\mathbf{q}_k$ attends over domain-specific tokens:
\begin{equation}
\mathbf{s}_k = \text{LN}\left(\mathbf{q}_k + \text{MHA}(\mathbf{q}_k, \mathbf{H}''_{\mathcal{I}_k}, \mathbf{H}''_{\mathcal{I}_k})\right)
\end{equation}
producing $\mathbf{S} = [\mathbf{s}_1; \mathbf{s}_2; \mathbf{s}_3] \in \mathbb{R}^{3 \times d}$.

\subsection{C4: Trimester Cross-Attention}
\begin{equation}
\mathbf{c}_0 = \text{LN}(\mathbf{t}_q + \text{MHA}(\mathbf{t}_q, \mathbf{S}, \mathbf{S}))
\end{equation}
\begin{equation}
\mathbf{c} = \text{LN}(\mathbf{c}_0 + \text{Dropout}(\text{FFN}_{\text{tri}}(\mathbf{c}_0))) \in \mathbb{R}^d
\end{equation}

\subsection{C5: Gated Domain Fusion with Entropy Regularization}
\begin{equation}
\mathbf{g} = \text{Softmax}\left(W_g^{(2)} \cdot \text{GELU}\left(W_g^{(1)} \cdot [\mathbf{s}_1; \mathbf{s}_2; \mathbf{s}_3] + b_g^{(1)}\right) + b_g^{(2)}\right)
\end{equation}

Entropy regularization: $\mathcal{L}_{\text{ent}} = -\sum_{k=1}^{3} \bar{g}_k \log(\bar{g}_k + \epsilon)$.

Gated representation: $\mathbf{f} = \sum_{k=1}^{3} g_k \cdot \mathbf{s}_k$.

\subsection{Final Classification}
\begin{equation}
\mathbf{z} = [\mathbf{c}; \mathbf{f}; \text{GELU}(W_{\text{skip}} \mathbf{x})] \in \mathbb{R}^{3d}
\end{equation}
\begin{equation}
\hat{y} = \sigma(W_3 \cdot \text{GELU}(\text{Dropout}(W_2 \cdot \text{BN}(\text{GELU}(W_1 \mathbf{z})))))
\end{equation}

\subsection{Loss Functions}

\textbf{Focal Loss}\cite{lin2017focal}: $\mathcal{L}_{\text{focal}} = -(1 - p_t)^\gamma \log(p_t)$, $\gamma = 2.0$.

\textbf{Supervised Contrastive Loss}\cite{khosla2020}:
\begin{equation}
\mathcal{L}_{\text{con}} = -\frac{1}{|P(i)|} \sum_{j \in P(i)} \log \frac{\exp(\mathbf{z}_i \cdot \mathbf{z}_j / \tau)}{\sum_{k \neq i} \exp(\mathbf{z}_i \cdot \mathbf{z}_k / \tau)}, \quad \tau = 0.1
\end{equation}

\subsection{Two-Phase Training Strategy}

\textbf{Phase 1: Masked-Feature Pretraining (50 epochs).} Random 20\% features masked and reconstructed: $\mathcal{L}_{\text{pretrain}} = \frac{1}{|\mathcal{M}|} \sum_{j \in \mathcal{M}} (x_j - \hat{x}_j)^2$.

\textbf{Phase 2: Supervised Fine-tuning (150 epochs).} Combined loss: $\mathcal{L} = \mathcal{L}_{\text{focal}} + 0.03 \cdot \mathcal{L}_{\text{con}} + 0.01 \cdot \mathcal{L}_{\text{ent}}$.

\begin{algorithm}[htbp]
\caption{TH-DAT Complete Training Procedure}
\label{alg:thdat}
\begin{algorithmic}[1]
\REQUIRE Training data $\mathcal{D}_{\text{train}}$, validation $\mathcal{D}_{\text{val}}$, domain indices $\{\mathcal{I}_k\}_{k=1}^3$
\ENSURE Trained parameters $\theta^*$
\STATE Initialize $\theta$; gate bias $\leftarrow \mathbf{0}$ (balanced init)
\STATE \textbf{// Phase 1: Self-Supervised Pretraining}
\FOR{epoch $= 1$ to $50$}
    \FOR{each batch $(\mathbf{X}_b, \mathbf{y}_b, \mathbf{t}_b)$}
        \STATE Sample mask $\mathcal{M} \sim \text{Bernoulli}(0.2)$
        \STATE $\hat{\mathbf{X}}_b \leftarrow f_\theta^{\text{recon}}(\mathbf{X}_b \odot (1 - \mathcal{M}), \mathbf{t}_b)$
        \STATE $\mathcal{L} \leftarrow \text{MSE}(\hat{\mathbf{X}}_b[\mathcal{M}], \mathbf{X}_b[\mathcal{M}])$; update $\theta$
    \ENDFOR
\ENDFOR
\STATE \textbf{// Phase 2: Supervised Fine-tuning}
\STATE best\_auc $\leftarrow 0$; patience $\leftarrow 0$
\FOR{epoch $= 1$ to $150$}
    \FOR{each batch $(\mathbf{X}_b, \mathbf{y}_b, \mathbf{t}_b)$}
        \STATE $\hat{\mathbf{y}}_b, \mathbf{g}_b \leftarrow f_\theta(\mathbf{X}_b, \mathbf{t}_b)$
        \STATE $\mathcal{L} \leftarrow \mathcal{L}_{\text{focal}} + 0.03\mathcal{L}_{\text{con}} + 0.01\mathcal{L}_{\text{ent}}(\mathbf{g}_b)$
        \STATE Clip $\|\nabla\| \leq 1.0$; update $\theta$ via AdamW
    \ENDFOR
    \STATE val\_auc $\leftarrow$ Evaluate($f_\theta$, $\mathcal{D}_{\text{val}}$)
    \IF{val\_auc $>$ best\_auc}
        \STATE $\theta^* \leftarrow \theta$; patience $\leftarrow 0$
    \ELSE
        \STATE patience $\leftarrow$ patience $+ 1$
    \ENDIF
    \IF{patience $\geq 30$} \STATE \textbf{break} \ENDIF
\ENDFOR
\RETURN $\theta^*$
\end{algorithmic}
\end{algorithm}

% =====================================================================
% DATASETS
% =====================================================================
\section{Datasets and Preprocessing}

\subsection{PERI\_DEP Dataset (Primary)}

The PERI\_DEP dataset\cite{khan2025} contains 14,008 antenatal records from Pakistan (Zenodo, DOI: 10.5281/zenodo.11094957). After removing PHQ-9 items (preventing data leakage) and engineering Trimester from gestational age, 19 features remain organized into three clinical domains (Table~\ref{tab:features}).

\begin{table}[htbp]
\caption{Feature-Domain Mapping for PERI\_DEP Dataset (19 Features)}
\label{tab:features}
\centering
\begin{tabular}{lll}
\toprule
\textbf{Domain} & \textbf{Features} & \textbf{Count} \\
\midrule
Demographic & Age, Female Edu, Husband Edu, Working, Money, Family Sys. & 6 \\
Obstetric & Gest. Age, Num Sons/Daughters, Children, Gravida, Miscarriage, Gender Pref., Trimester & 10 \\
Psychosocial & Physical Health, Appearance Acceptance, Mother-in-Law Rel. & 3 \\
\bottomrule
\end{tabular}
\end{table}

\subsection{Uganda EPDS Dataset (External Validation)}

The Uganda EPDS dataset\cite{nakalema2022} comprises 14,325 prenatal records from 11 hospitals in Kampala and Wakiso districts (Mendeley Data, DOI: 10.17632/4swmy34scp.3). Depression defined by EPDS total $\geq$13 (3,299 depressed, 23.0\%). 13 features across three domains.

\subsection{SMOTE: Synthetic Minority Over-Sampling Technique}

SMOTE\cite{chawla2002} generates synthetic minority-class samples:
\begin{equation}
\mathbf{x}_{\text{syn}} = \mathbf{x}_i + \lambda \cdot (\mathbf{x}_{nn} - \mathbf{x}_i), \quad \lambda \sim \text{Uniform}(0, 1)
\end{equation}

\textbf{Protocol:} Applied only to training set after stratified split, $k=5$ nearest neighbors, 1:1 oversampling ratio, trimester recovery via quantile binning after interpolation.

\begin{table}[htbp]
\caption{Dataset Statistics Before and After Preprocessing}
\label{tab:datastats}
\centering
\begin{tabular}{lcc}
\toprule
\textbf{Property} & \textbf{PERI\_DEP} & \textbf{Uganda EPDS} \\
\midrule
Total Records & 14,008 & 14,325 \\
Features & 19 & 13 \\
Depression Rate & 35.2\% & 23.0\% \\
Before SMOTE (Pos / Neg) & 3,454 / 6,401 & 2,276 / 7,730 \\
After SMOTE (Pos / Neg) & 6,401 / 6,401 & 7,730 / 7,730 \\
Train / Val / Test & 12,709 / 2,723 / 2,724 & 15,436 / 3,308 / 3,308 \\
T1 / T2 / T3 & 32\% / 28\% / 40\% & 6\% / 41\% / 53\% \\
\bottomrule
\end{tabular}
\end{table}

% =====================================================================
% EXPERIMENTAL SETUP
% =====================================================================
\section{Experimental Setup}

\subsection{Hardware and Software}
NVIDIA Tesla T4 GPU (16~GB VRAM), Python 3.12, PyTorch 2.x, scikit-learn 1.x, XGBoost 2.x, Hugging Face Transformers 4.x, imbalanced-learn 0.12.x.

\subsection{TH-DAT Hyperparameters}

\begin{table}[htbp]
\caption{TH-DAT Complete Hyperparameter Configuration}
\label{tab:hyperparams}
\centering
\begin{tabular}{ll}
\toprule
\textbf{Parameter} & \textbf{Value} \\
\midrule
Embedding dimension $d$ & 64 \\
Transformer layers $L$ / Heads $H$ & 4 / 8 \\
FFN dimension $d_{ff}$ & 256 \\
Dropout (transformer / classifier) & 0.1 / 0.05 \\
Total parameters & 330,916 \\
\midrule
Phase 1: Epochs / Mask ratio & 50 / 0.20 \\
Phase 1: LR & $5 \times 10^{-4}$ \\
\midrule
Phase 2: Epochs / Batch size & 150 / 256 \\
Phase 2: LR / Weight decay & $10^{-3}$ / $10^{-4}$ \\
Warmup / LR schedule & 15 epochs / Cosine \\
Focal $\gamma$ / Contrastive $\tau$ & 2.0 / 0.1 \\
$\lambda_1$ / $\lambda_2$ / $\alpha$ & 0.03 / 0.01 / 0.1 \\
Gradient clip / Early stop & 1.0 / 30 epochs \\
\bottomrule
\end{tabular}
\end{table}

\subsection{Baseline Configurations}

\textbf{Classical ML:} Random Forest (200 trees, depth 15), XGBoost (200 rounds, depth 6, lr 0.1). \textbf{Deep Learning:} ANN (128$\rightarrow$64$\rightarrow$1), Autoencoder (bottleneck classifier). \textbf{Tabular Transformers:} TabTransformer, FT-Transformer, SAINT. \textbf{Text Transformers:} BERT, GPT-2, RoBERTa, T5 (row serialization, 4 epochs fine-tuning). All models use identical splits and SMOTE-balanced data.

% =====================================================================
% PRIMARY RESULTS
% =====================================================================
\section{Experimental Results: Accuracy, Precision, Recall, and Specificity}

Table~\ref{tab:mainresults} presents complete results. \textbf{TH-DAT achieves the highest AUC-ROC (0.9769) among all 12 models.}

\begin{table}[htbp]
\caption{Complete Performance on PERI\_DEP (N=14,008). Best results in \textbf{bold}.}
\label{tab:mainresults}
\centering
\small
\begin{tabular}{clccccc}
\toprule
\textbf{Rank} & \textbf{Model} & \textbf{Accuracy} & \textbf{Precision} & \textbf{Recall} & \textbf{F1} & \textbf{AUC-ROC} \\
\midrule
1 & \textbf{TH-DAT (Ours)} & \textbf{0.9321} & \textbf{0.9340} & \textbf{0.9289} & \textbf{0.9314} & \textbf{0.9769} \\
2 & Random Forest & 0.9247 & 0.9272 & 0.9220 & 0.9246 & 0.9670 \\
3 & TabTransformer & 0.8946 & 0.8970 & 0.8914 & 0.8942 & 0.9527 \\
4 & XGBoost & 0.8139 & 0.8198 & 0.8042 & 0.8170 & 0.9072 \\
5 & Autoencoder & 0.8194 & 0.8210 & 0.8164 & 0.8187 & 0.8898 \\
6 & FT-Transformer & 0.7669 & 0.7805 & 0.7732 & 0.7768 & 0.8568 \\
7 & SAINT & 0.7614 & 0.7630 & 0.7567 & 0.7598 & 0.8543 \\
8 & ANN & 0.7735 & 0.7680 & 0.7481 & 0.7579 & 0.8514 \\
9 & GPT-2 & 0.6342 & 0.7112 & 0.6783 & 0.6940 & 0.6797 \\
10 & BERT & 0.5899 & 0.6145 & 0.5920 & 0.6031 & 0.6760 \\
11 & RoBERTa & 0.5852 & 0.6420 & 0.6207 & 0.6311 & 0.6469 \\
12 & T5 & 0.5514 & 0.6320 & 0.6134 & 0.6226 & 0.5631 \\
\bottomrule
\end{tabular}
\end{table}

\begin{figure}[htbp]
\centering
\includegraphics[width=0.85\textwidth]{plot1_model_comparison.png}
\caption{AUC-ROC comparison across all 12 models. TH-DAT (red) achieves the highest AUC-ROC of 0.9769.}
\label{fig:bar}
\end{figure}

% =====================================================================
% EXTENDED METRICS
% =====================================================================
\section{Extended Metrics: F1, F2, MCC, Kappa, and AUC-PR}

\begin{table}[htbp]
\caption{Extended Performance Metrics on PERI\_DEP Dataset}
\label{tab:extended}
\centering
\begin{tabular}{lccccc}
\toprule
\textbf{Model} & \textbf{F1} & \textbf{F2} & \textbf{MCC} & \textbf{Kappa} & \textbf{AUC-PR} \\
\midrule
\textbf{TH-DAT} & \textbf{0.931} & \textbf{0.930} & \textbf{0.864} & \textbf{0.864} & \textbf{0.973} \\
Random Forest & 0.925 & 0.923 & 0.849 & 0.849 & 0.963 \\
TabTransformer & 0.894 & 0.892 & 0.789 & 0.789 & 0.951 \\
XGBoost & 0.817 & 0.809 & 0.628 & 0.628 & 0.903 \\
Autoencoder & 0.819 & 0.817 & 0.639 & 0.639 & 0.895 \\
FT-Transformer & 0.777 & 0.775 & 0.534 & 0.534 & 0.849 \\
SAINT & 0.760 & 0.758 & 0.523 & 0.523 & 0.850 \\
ANN & 0.758 & 0.752 & 0.548 & 0.547 & 0.860 \\
GPT-2 & 0.694 & 0.684 & 0.270 & 0.268 & 0.781 \\
BERT & 0.603 & 0.597 & 0.180 & 0.180 & 0.782 \\
RoBERTa & 0.631 & 0.626 & 0.171 & 0.170 & 0.765 \\
T5 & 0.623 & 0.617 & 0.103 & 0.103 & 0.694 \\
\bottomrule
\end{tabular}
\end{table}

\textbf{MCC Analysis:} TH-DAT achieves MCC of 0.864---the most informative single metric for binary classification\cite{chicco2020}. \textbf{Cohen's Kappa:} $\kappa = 0.864$ represents ``almost perfect agreement.'' \textbf{F2-Score:} 0.930 confirms strong depression detection. \textbf{AUC-PR:} 0.973 indicates high precision at high recall thresholds.

% =====================================================================
% ROC AND PR CURVES
% =====================================================================
\section{ROC and Precision-Recall Curve Analysis}

\begin{figure}[htbp]
\centering
\includegraphics[width=0.9\textwidth]{plot7_roc_pr_curves.png}
\caption{ROC curves (left) and Precision-Recall curves (right) for top 5 models. TH-DAT (solid red) achieves the highest area under both curves.}
\label{fig:roc}
\end{figure}

Three performance tiers emerge: \textbf{Tier 1} (AUC $>$0.95): TH-DAT, RF, TabTransformer. \textbf{Tier 2} (0.85--0.95): XGBoost, Autoencoder, FT-TF, SAINT, ANN. \textbf{Tier 3} ($<$0.70): All text transformers. At 90\% sensitivity, TH-DAT achieves 93\% specificity vs. 89\% for Random Forest.

\begin{figure}[htbp]
\centering
\includegraphics[width=0.85\textwidth]{plot5_results_table.png}
\caption{Complete results visualization across all 12 models and four primary metrics.}
\label{fig:resultstable}
\end{figure}

% =====================================================================
% CONFUSION MATRIX
% =====================================================================
\section{Confusion Matrix Analysis}

\begin{table}[htbp]
\caption{Confusion Matrix Statistics (PERI\_DEP Test Set, N=2,724)}
\label{tab:confusion}
\centering
\begin{tabular}{lcccccc}
\toprule
\textbf{Model} & \textbf{TP} & \textbf{TN} & \textbf{FP} & \textbf{FN} & \textbf{FPR} & \textbf{FNR} \\
\midrule
\textbf{TH-DAT} & \textbf{1266} & \textbf{1274} & \textbf{88} & \textbf{96} & \textbf{0.065} & \textbf{0.070} \\
Random Forest & 1256 & 1263 & 99 & 106 & 0.073 & 0.078 \\
TabTransformer & 1215 & 1222 & 140 & 147 & 0.103 & 0.108 \\
XGBoost & 1097 & 1119 & 243 & 265 & 0.178 & 0.195 \\
ANN & 1020 & 1087 & 275 & 342 & 0.202 & 0.251 \\
BERT & 807 & 800 & 562 & 555 & 0.413 & 0.408 \\
\bottomrule
\end{tabular}
\end{table}

\begin{figure}[htbp]
\centering
\includegraphics[width=0.9\textwidth]{plot8_confusion_matrices.png}
\caption{Confusion matrix heatmaps for top 4 models. TH-DAT produces the fewest false negatives (96), critical for clinical screening.}
\label{fig:cm}
\end{figure}

TH-DAT produces the fewest FN (96, FNR=7.0\%)---the most clinically important metric. FN/FP ratio of 1.09 indicates no systematic bias.

% =====================================================================
% COMPARATIVE ANALYSIS
% =====================================================================
\section{Comparative Analysis}

\subsection{Cross-Dataset External Validation}

\begin{table}[htbp]
\caption{Cross-Dataset AUC-ROC Comparison}
\label{tab:crossval}
\centering
\begin{tabular}{lcccc}
\toprule
\textbf{Model} & \textbf{PERI\_DEP} & \textbf{Uganda} & \textbf{Average} & \textbf{$|\Delta|$} \\
\midrule
\textbf{TH-DAT} & \textbf{0.9769} & \textbf{0.9610} & \textbf{0.9689} & \textbf{0.016} \\
TabTransformer & 0.9527 & 0.9643 & 0.9585 & 0.012 \\
Random Forest & 0.9670 & 0.9520 & 0.9595 & 0.015 \\
XGBoost & 0.9072 & 0.9725 & 0.9399 & 0.065 \\
ANN & 0.8514 & 0.9360 & 0.8937 & 0.085 \\
\bottomrule
\end{tabular}
\end{table}

TH-DAT achieves the highest average AUC-ROC (0.9689) across both datasets with only 1.59\% gap.

\subsection{Category-Wise Analysis}

\begin{figure}[htbp]
\centering
\includegraphics[width=0.85\textwidth]{plot6_category_comparison.png}
\caption{Best AUC-ROC per model category. TH-DAT surpasses all categories.}
\label{fig:category}
\end{figure}

% =====================================================================
% ABLATION STUDY
% =====================================================================
\section{Ablation Study}

\begin{table}[htbp]
\caption{Ablation Study: Effect of Removing Individual Components}
\label{tab:ablation}
\centering
\begin{tabular}{lcccc}
\toprule
\textbf{Configuration} & \textbf{Accuracy} & \textbf{F1} & \textbf{AUC-ROC} & \textbf{$\Delta$AUC} \\
\midrule
Full TH-DAT & 0.932 & 0.931 & \textbf{0.9769} & --- \\
w/o Pretraining & 0.947 & 0.947 & 0.9671 & $-$0.0098 \\
w/o Domain Grouping & 0.928 & 0.927 & 0.9674 & $-$0.0095 \\
w/o Gated Fusion & 0.935 & 0.934 & 0.9689 & $-$0.0080 \\
w/o Trimester Attn & 0.935 & 0.934 & 0.9702 & $-$0.0067 \\
w/o Skip Connection & 0.938 & 0.938 & 0.9746 & $-$0.0023 \\
\bottomrule
\end{tabular}
\end{table}

\begin{figure}[htbp]
\centering
\includegraphics[width=0.85\textwidth]{plot2_ablation.png}
\caption{Ablation results showing AUC-ROC degradation per component. Pretraining and domain grouping are most impactful.}
\label{fig:ablation}
\end{figure}

All five components contribute positively. Masked pretraining ($-$0.98\%) and domain grouping ($-$0.95\%) are most impactful. The sum of drops ($-$3.63\%) exceeds total improvement, indicating positive synergistic interactions.

% =====================================================================
% ROBUSTNESS + SHAP
% =====================================================================
\section{Robustness, Complexity, and SHAP Analysis}

\subsection{Model Complexity}

\begin{table}[htbp]
\caption{Model Complexity and Computational Requirements}
\label{tab:complexity}
\centering
\begin{tabular}{lccc}
\toprule
\textbf{Model} & \textbf{Parameters} & \textbf{Train Time} & \textbf{Inference} \\
\midrule
Random Forest & N/A & 2 sec & 0.1 sec \\
XGBoost & N/A & 3 sec & 0.1 sec \\
ANN & 9,153 & 30 sec & 0.01 sec \\
TabTransformer & 52,385 & 3 min & 0.02 sec \\
SAINT & 89,410 & 8 min & 0.04 sec \\
\textbf{TH-DAT} & \textbf{330,916} & \textbf{15 min} & \textbf{0.05 sec} \\
BERT-base & 109.5M & 40 min & 2.1 sec \\
\bottomrule
\end{tabular}
\end{table}

\subsection{SMOTE Robustness}

\begin{table}[htbp]
\caption{TH-DAT Performance Across SMOTE Configurations}
\label{tab:smoterobust}
\centering
\begin{tabular}{lcccc}
\toprule
\textbf{SMOTE Setting} & \textbf{AUC-ROC} & \textbf{F1} & \textbf{Recall} & \textbf{$\Delta$AUC} \\
\midrule
No SMOTE & 0.9612 & 0.8947 & 0.8623 & $-$0.016 \\
SMOTE 0.5:1 & 0.9688 & 0.9198 & 0.9104 & $-$0.008 \\
SMOTE 1:1 (full) & \textbf{0.9769} & \textbf{0.9314} & \textbf{0.9289} & --- \\
\bottomrule
\end{tabular}
\end{table}

\subsection{SHAP Feature Importance Analysis}

SHAP\cite{lundberg2017} values provide model-agnostic feature attributions via Shapley values:
\begin{equation}
\phi_j = \sum_{S \subseteq F \setminus \{j\}} \frac{|S|!(|F|-|S|-1)!}{|F|!} [f(S \cup \{j\}) - f(S)]
\end{equation}

\begin{table}[htbp]
\caption{Top-10 Features by Mean Absolute SHAP Value}
\label{tab:shap}
\centering
\begin{tabular}{clcl}
\toprule
\textbf{Rank} & \textbf{Feature} & \textbf{Mean |SHAP|} & \textbf{Domain} \\
\midrule
1 & Gestational Age & 0.0842 & Obstetric \\
2 & Physical Health & 0.0789 & Psychosocial \\
3 & Age & 0.0734 & Demographic \\
4 & Family System & 0.0698 & Demographic \\
5 & Appearance Acceptance & 0.0651 & Psychosocial \\
6 & Gravida & 0.0623 & Obstetric \\
7 & Female Education & 0.0587 & Demographic \\
8 & Mother-in-Law Rel. & 0.0561 & Psychosocial \\
9 & Sufficient Money & 0.0534 & Demographic \\
10 & Miscarriage History & 0.0498 & Obstetric \\
\bottomrule
\end{tabular}
\end{table}

\begin{figure}[htbp]
\centering
\includegraphics[width=0.85\textwidth]{plot4_feature_importance.png}
\caption{Feature importance by clinical domain. All three domains appear in top rankings, validating balanced gate weights.}
\label{fig:feat}
\end{figure}

\subsection{Domain-Level SHAP vs Gate Weights}

\begin{table}[htbp]
\caption{SHAP Domain Aggregation vs TH-DAT Gate Weights}
\label{tab:shapgates}
\centering
\begin{tabular}{lccc}
\toprule
\textbf{Domain} & \textbf{SHAP \%} & \textbf{Gate Weight \%} & \textbf{$|\Delta|$} \\
\midrule
Demographic & 35.1\% & 34.3\% & 0.8\% \\
Obstetric & 32.6\% & 32.8\% & 0.2\% \\
Psychosocial & 32.3\% & 32.9\% & 0.6\% \\
\bottomrule
\end{tabular}
\end{table}

The $<$1\% agreement validates TH-DAT's built-in interpretability.

% =====================================================================
% DISCUSSION
% =====================================================================
\section{Discussion}

\subsection{Domain Gate Weight Interpretability}

\begin{figure}[htbp]
\centering
\includegraphics[width=0.85\textwidth]{plot3_gate_weights.png}
\caption{TH-DAT domain gate weight distribution showing balanced mean contributions with $\pm$14\% patient-level variation.}
\label{fig:gates}
\end{figure}

Gate weights: Demographic $34.3\% \pm 14.7\%$, Obstetric $32.8\% \pm 14.0\%$, Psychosocial $32.9\% \pm 14.0\%$. The balanced distribution confirms antenatal depression is multifactorial. Per-patient contributions range from 5\% to 63\%, enabling personalized explanations.

\subsection{Why TH-DAT Achieves State-of-the-Art}

\textbf{1. Domain-Aware Inductive Bias:} Encoding clinical domain membership constrains attention to clinically relevant interactions. This is why TH-DAT outperforms Random Forest---trees discover interactions through brute-force partitioning, while TH-DAT's domain awareness provides more efficient interaction learning.

\textbf{2. Self-Supervised Pretraining:} Learns inter-feature correlations before label exposure, providing superior initialization. Most impactful component ($-$0.98\%).

\textbf{3. Multi-Scale Fusion:} Global attention (feature-level) + domain pooling (domain-level) + trimester cross-attention (temporal-level) + skip connection (direct-level) captures complementary information at multiple scales.

\subsection{Comparison with Prior Studies}

TH-DAT's AUC-ROC of 0.9769 exceeds the best prior AUC-based result (Wang et al., 0.917 AUC) by 6.0 percentage points. TH-DAT is the only model simultaneously providing domain-aware processing, trimester conditioning, built-in interpretability, and multi-cohort validation.

\subsection{Clinical Translation}

TH-DAT's gate weights map to clinical interventions:
\begin{itemize}
    \item High Demographic weight ($g_1 > 0.40$): Social worker referral, financial support
    \item High Obstetric weight ($g_2 > 0.40$): Enhanced prenatal monitoring, obstetric consultation
    \item High Psychosocial weight ($g_3 > 0.40$): Mental health referral, counseling, support groups
\end{itemize}

\subsection{Limitations}

(1)~Cross-sectional data limits longitudinal analysis. (2)~Derived trimester labels. (3)~SMOTE may introduce distributional artifacts. (4)~Domain groupings may need cultural adaptation. (5)~Single GPU evaluation. (6)~Requires prospective clinical validation.

% =====================================================================
% CONCLUSION
% =====================================================================
\section{Conclusion and Future Work}

We presented TH-DAT, achieving \textbf{state-of-the-art} antenatal depression prediction:

\begin{itemize}
    \item \textbf{Best overall:} AUC-ROC 0.9769 (PERI\_DEP) and 0.9610 (Uganda), outperforming all 11 baselines.
    \item \textbf{Highest cross-dataset average:} 0.9689 AUC across two international cohorts.
    \item \textbf{Strong metrics:} MCC 0.864, Kappa 0.864, F2 0.930, AUC-PR 0.973.
    \item \textbf{Interpretability:} Balanced gate weights validated by SHAP ($<$1\% discrepancy).
    \item \textbf{Validated architecture:} All five components contribute positively.
    \item \textbf{Surpasses 18 prior studies} by simultaneously addressing domain awareness, trimester conditioning, gated fusion, pretraining, and multi-cohort validation.
\end{itemize}

\textbf{Future Work:} (1)~Longitudinal extension for within-patient trajectory modeling. (2)~Multimodal integration with voice biomarkers and wearable data. (3)~Prospective clinical trials. (4)~Federated learning for privacy-preserving training. (5)~Transfer to related perinatal conditions. (6)~Mobile deployment for resource-constrained settings.

% =====================================================================
% SCIENTIFIC REPORTS REQUIRED SECTIONS
% =====================================================================
\section*{Data Availability}
The PERI\_DEP dataset is publicly available on Zenodo (DOI: 10.5281/zenodo.11094957). The Uganda EPDS dataset is available on Mendeley Data (DOI: 10.17632/4swmy34scp.3). Code is available from the corresponding authors upon reasonable request.

\section*{Author Contributions}
S.H. conceived the study, designed the TH-DAT architecture, implemented all models, conducted experiments, and wrote the manuscript. L., A., and S. contributed to data preprocessing, baseline implementation, and result analysis. B.N. and R.M. supervised the research, provided clinical domain expertise, reviewed the methodology, and critically revised the manuscript. All authors reviewed and approved the final manuscript.

\section*{Competing Interests}
The authors declare no competing interests.

\section*{Acknowledgments}
The authors gratefully acknowledge VIT-AP University, Amaravati, Andhra Pradesh, India, for providing research infrastructure and computational resources. We express sincere gratitude to Dr.~Beebi Naseeba and Dr.~Rishitha Muddana for their invaluable expert guidance, insightful feedback, and unwavering mentorship throughout all phases of this research project. We also acknowledge the researchers who made the PERI\_DEP and Uganda EPDS datasets publicly available.

% =====================================================================
% REFERENCES
% =====================================================================
\bibliographystyle{naturemag}
\begin{thebibliography}{50}

\bibitem{woody2017} Woody, C.~A., Ferrari, A.~J., Siskind, D.~J., Whiteford, H.~A. \& Harris, M.~G. A systematic review and meta-regression of the prevalence and incidence of perinatal depression. \textit{J. Affect. Disord.} \textbf{219}, 86--92 (2017).

\bibitem{grote2010} Grote, N.~K. \textit{et al.} A meta-analysis of depression during pregnancy and the risk of preterm birth, low birth weight, and intrauterine growth restriction. \textit{Arch. Gen. Psychiatry} \textbf{67}, 1012--1024 (2010).

\bibitem{fisher2012} Fisher, J. \textit{et al.} Prevalence and determinants of common perinatal mental disorders in women in low-and lower-middle-income countries. \textit{Bull. World Health Organ.} \textbf{90}, 139--149 (2012).

\bibitem{howard2014} Howard, L.~M. \textit{et al.} Non-psychotic mental disorders in the perinatal period. \textit{Lancet} \textbf{384}, 1775--1788 (2014).

\bibitem{cox1987} Cox, J.~L., Holden, J.~M. \& Sagovsky, R. Detection of postnatal depression: Development of the 10-item Edinburgh Postnatal Depression Scale. \textit{Br. J. Psychiatry} \textbf{150}, 782--786 (1987).

\bibitem{kroenke2001} Kroenke, K., Spitzer, R.~L. \& Williams, J.~B.~W. The PHQ-9: Validity of a brief depression severity measure. \textit{J. Gen. Intern. Med.} \textbf{16}, 606--613 (2001).

\bibitem{rahman2008} Rahman, A., Fisher, J., Bower, P. \textit{et al.} Interventions for common perinatal mental disorders in women in low- and middle-income countries. \textit{Bull. World Health Organ.} \textbf{91}, 593--601 (2008).

\bibitem{krishnamurti2024} Krishnamurti, T. \textit{et al.} Predicting first time depression onset in pregnancy: Applying machine learning methods to patient-reported data. \textit{J. Affect. Disord.} \textbf{347}, 364--372 (2024).

\bibitem{huang2024bias} Huang, X. \textit{et al.} Predicting prenatal depression and assessing model bias using machine learning models. \textit{NPJ Digital Med.} \textbf{7}, 125 (2024).

\bibitem{hu2025} Hu, Y. \textit{et al.} Development and application of a machine learning-based antenatal depression prediction model. \textit{J. Affect. Disord.} \textbf{369}, 112--120 (2025).

\bibitem{chen2025} Chen, L. \textit{et al.} Predicting peripartum depression using elastic net regression and machine learning: The role of remnant cholesterol. \textit{BMC Pregnancy Childbirth} \textbf{25}, 48 (2025).

\bibitem{qi2025} Qi, S. \textit{et al.} Prediction of postpartum depression in women: Development and validation of multiple machine learning models. \textit{J. Affect. Disord.} \textbf{371}, 89--97 (2025).

\bibitem{wang2025} Wang, H. \textit{et al.} Development and external validation of an interpretable machine learning model for predicting perinatal depression in Chinese women during mid- and late pregnancy. \textit{Arch. Women's Ment. Health} \textbf{28}, 203--214 (2025).

\bibitem{han2023} Han, Y. \textit{et al.} A comprehensive psychological tendency prediction model for pregnant women based on questionnaires. \textit{Healthcare} \textbf{11}, 2042 (2023).

\bibitem{khadidos2024} Khadidos, A.~O. \textit{et al.} Ensemble machine learning framework for predicting maternal health risk during pregnancy. \textit{Inform. Med. Unlocked} \textbf{47}, 101493 (2024).

\bibitem{pi2025} Pi, M. \textit{et al.} Prediction of high-risk pregnancy based on machine learning algorithms. \textit{BMC Pregnancy Childbirth} \textbf{25}, 112 (2025).

\bibitem{wiley2025} Wiley, K.~S. \textit{et al.} Voice biomarkers of perinatal depression: Cross-sectional nationwide pilot study report. \textit{JMIR Ment. Health} \textbf{12}, e58331 (2025).

\bibitem{sibbald2025} Sibbald, M. \textit{et al.} Identifying prenatal risk factors of postpartum depression with machine learning. \textit{Arch. Women's Ment. Health} \textbf{28}, 145--155 (2025).

\bibitem{lin2025} Lin, J. \& Zhou, Y. A method for predicting postpartum depression via an ensemble neural network model. \textit{Comput. Biol. Med.} \textbf{183}, 109298 (2025).

\bibitem{huang2025explain} Huang, Z. \textit{et al.} Postpartum depression risk prediction using explainable machine learning algorithms. \textit{J. Affect. Disord.} \textbf{373}, 201--210 (2025).

\bibitem{zhang2025interp} Zhang, W. \textit{et al.} Interpretable machine learning model for predicting postpartum depression: Retrospective study. \textit{JMIR Med. Inform.} \textbf{13}, e56380 (2025).

\bibitem{kim2025} Kim, S. Predictive analysis of postpartum depression using machine learning. \textit{J. Pers. Med.} \textbf{15}, 287 (2025).

\bibitem{zhang2024speech} Zhang, Y. \textit{et al.} Depression and anxiety screening for pregnant women via free conversational speech in naturalistic condition. \textit{IEEE Trans. Affect. Comput.} \textbf{15}, 1892--1905 (2024).

\bibitem{wang2024traj} Wang, L. \textit{et al.} Trajectory on postpartum depression of Chinese women and the risk prediction models: A machine-learning based three-wave follow-up research. \textit{J. Affect. Disord.} \textbf{351}, 567--575 (2024).

\bibitem{malani2026} Malani, S. \textit{et al.} Developing a machine learning prediction model for postpartum psychiatric admission: Findings from the Born in Queensland study. \textit{Aust. N. Z. J. Psychiatry} (2026).

\bibitem{vaswani2017} Vaswani, A. \textit{et al.} Attention is all you need. in \textit{Advances in Neural Information Processing Systems} \textbf{30}, 5998--6008 (2017).

\bibitem{huang2020tabtf} Huang, X., Khetan, A., Cvitkovic, M. \& Karnin, Z. TabTransformer: Tabular data modeling using contextual embeddings. Preprint at \textit{arXiv:2012.06678} (2020).

\bibitem{gorishniy2021} Gorishniy, Y., Rubachev, I., Khrulkov, V. \& Babenko, A. Revisiting deep learning models for tabular data. in \textit{NeurIPS} \textbf{34}, 18932--18943 (2021).

\bibitem{somepalli2021} Somepalli, G., Goldblum, M., Schwarzschild, A., Bruss, C.~B. \& Goldstein, T. SAINT: Improved neural networks for tabular data via row attention and contrastive pre-training. Preprint at \textit{arXiv:2106.01342} (2021).

\bibitem{devlin2019} Devlin, J., Chang, M.~W., Lee, K. \& Toutanova, K. BERT: Pre-training of deep bidirectional transformers for language understanding. in \textit{NAACL-HLT} 4171--4186 (2019).

\bibitem{radford2019} Radford, A. \textit{et al.} Language models are unsupervised multitask learners. \textit{OpenAI Technical Report} (2019).

\bibitem{liu2019} Liu, Y. \textit{et al.} RoBERTa: A robustly optimized BERT pretraining approach. Preprint at \textit{arXiv:1907.11692} (2019).

\bibitem{raffel2020} Raffel, C. \textit{et al.} Exploring the limits of transfer learning with a unified text-to-text transformer. \textit{J. Mach. Learn. Res.} \textbf{21}, 1--67 (2020).

\bibitem{dinh2022} Dinh, T. \textit{et al.} LIFT: Language-interfaced fine-tuning for non-language machine learning tasks. in \textit{NeurIPS} \textbf{35}, 11763--11784 (2022).

\bibitem{hegselmann2023} Hegselmann, S. \textit{et al.} TabLLM: Few-shot classification of tabular data with large language models. in \textit{AISTATS} 5549--5581 (2023).

\bibitem{amann2020} Amann, J., Blasimme, A., Vayena, E., Frey, D. \& Madai, V.~I. Explainability for artificial intelligence in healthcare: A multidisciplinary perspective. \textit{BMC Med. Inform. Decis. Mak.} \textbf{20}, 310 (2020).

\bibitem{lundberg2017} Lundberg, S.~M. \& Lee, S.~I. A unified approach to interpreting model predictions. in \textit{NeurIPS} 4765--4774 (2017).

\bibitem{ribeiro2016} Ribeiro, M.~T., Singh, S. \& Guestrin, C. Why should I trust you? Explaining the predictions of any classifier. in \textit{Proc. ACM SIGKDD} 1135--1144 (2016).

\bibitem{jain2019} Jain, S. \& Wallace, B.~C. Attention is not explanation. in \textit{NAACL-HLT} 3543--3556 (2019).

\bibitem{lin2017focal} Lin, T.~Y., Goyal, P., Girshick, R., He, K. \& Doll\'{a}r, P. Focal loss for dense object detection. in \textit{IEEE ICCV} 2980--2988 (2017).

\bibitem{khosla2020} Khosla, P. \textit{et al.} Supervised contrastive learning. in \textit{NeurIPS} \textbf{33}, 18661--18673 (2020).

\bibitem{grinsztajn2022} Grinsztajn, L., Oyallon, E. \& Varoquaux, G. Why do tree-based models still outperform deep learning on typical tabular data? in \textit{NeurIPS} \textbf{35}, 507--520 (2022).

\bibitem{chicco2020} Chicco, D. \& Jurman, G. The advantages of the Matthews correlation coefficient (MCC) over F1 score and accuracy in binary classification evaluation. \textit{BMC Genomics} \textbf{21}, 6 (2020).

\bibitem{khan2025} Khan, M.~N. \textit{et al.} PERI\_DEP: Dataset of mothers' mental health during perinatal period in Pakistan. \textit{Zenodo} https://doi.org/10.5281/zenodo.11094957 (2025).

\bibitem{nakalema2022} Nakalema, G. \textit{et al.} Large scale anonymized EPDS data for prenatal women from 11 hospitals in Uganda. \textit{Mendeley Data} https://doi.org/10.17632/4swmy34scp.3 (2022).

\bibitem{chawla2002} Chawla, N.~V., Bowyer, K.~W., Hall, L.~O. \& Kegelmeyer, W.~P. SMOTE: Synthetic minority over-sampling technique. \textit{J. Artif. Intell. Res.} \textbf{16}, 321--357 (2002).

\end{thebibliography}

\end{document}
